\documentclass[12pt]{article}
\usepackage[T1]{fontenc}
\usepackage[utf8]{inputenc}
\usepackage{lmodern}
\usepackage{textcomp}
\usepackage{enumitem}
\usepackage{geometry}
\usepackage{graphicx}
\usepackage{amsmath, amssymb}
\geometry{margin=1in}
\begin{document}
\begin{center}
Biology - Graduate Level Assessment 4 \
Total Marks: 40 \
Answer all questions.
\end{center}

\section*{Multiple Choice (10 marks)}
\begin{enumerate}[label=\arabic*.]
\item Which of the following is the symplastic pathway for the movement of sucrose from the site of photosynthesis in mesophyll cells into the phloem?
\begin{enumerate}[label=(\alph*)]
    \item Bundle sheath, companion cell, phloem parenchyma, fibers
    \item Phloem parenchyma, bundle sheath, fibers, companion cell
    \item Bundle sheath, phloem parenchyma, companion cell, sieve tube
    \item Bundle sheath, fibers, companion cell, sieve tube
    \item Fibers, bundle sheath, companion cell, sieve tube
\end{enumerate}
\item What prevents the stomach from being digested by itsown secretions?
\begin{enumerate}[label=(\alph*)]
    \item The thick muscular wall of the stomach
    \item Antibacterial properties of stomach acid
    \item Gastric juice
    \item High pH levels in the stomach
    \item Presence of beneficial bacteria
\end{enumerate}
\item What is the physiological purpose of shivering?
\begin{enumerate}[label=(\alph*)]
    \item Shivering is a method the body uses to relieve pain.
    \item Shivering is a tactic to ward off predators by appearing agitated.
    \item Shivering is a way to burn excess calories
    \item Shivering is a response to fear or anxiety
    \item Shivering is an involuntary action to prepare muscles for physical activity.
\end{enumerate}
\item What are some of the cues used by animal in migrating?
\begin{enumerate}[label=(\alph*)]
    \item Taste, wind direction, humidity
    \item Atmospheric oxygen levels, plant growth, tidal patterns
    \item Infrared radiation, sound frequency, barometric pressure
    \item Moon phases, echo location, gravitational pull
    \item Sunlight intensity, water current, soil composition
\end{enumerate}
\item In certain abnormal conditions, the stomach does not secretehydrochloric acid. What effects might this haveon thedigestive process?
\begin{enumerate}[label=(\alph*)]
    \item Without HCl, more bacteria would be killed in the stomach.
    \item No HCl would lead to reduced mucus production in the stomach lining.
    \item The absence of HCl would cause the stomach to digest its own tissues.
    \item No HCl in the stomach would significantly increase the pH of the small intestine.
    \item The absence of HCl would speed up the digestive process.
\end{enumerate}
\end{enumerate}

\section*{True / False (10 marks)}
\textit{Answer True or False}
\begin{enumerate}[label=\arabic*.]
\setcounter{enumi}{5}
\item True or False: Does anterior laxity of the uninjured knee influence clinical outcomes of ACL reconstruction?
\item True or False: Are UK radiologists satisfied with the training and support received in suspected child abuse?
\item True or False: Must early postoperative oral intake be limited to laparoscopy?
\item True or False: Can the prognosis of polymyalgia rheumatica be predicted at disease onset?
\item True or False: Thrombosis prophylaxis in hospitalised medical patients: does prophylaxis in all patients make sense?
\end{enumerate}

\section*{Long Form Response (20 marks)}
\textit{Answer each question with a clear and complete explanation.}
\begin{enumerate}[label=\arabic*.]
\setcounter{enumi}{10}
\item Discuss the challenges of using human subjects in inheritance studies, considering genetic complexity, environmental factors, ethical dilemmas, and family dynamics.
\item Discuss the factors that promote linkage disequilibrium in populations, and analyze why random mating does not contribute to this phenomenon.
\end{enumerate}

\vfill
\noindent\textit{End of Paper}
\end{document}